\section{Implementation Surface}
\label{sec:receipts-implementation}

\subsection{Source Files}
\label{sec:receipts-implementation-source}

The receipts corpus is realized across fifteen source files spanning three organizational
layers. The M\&A claim receipt layer contains five signed JSON documents. The
\texttt{rec\_ebitda\_rework\_001.json} receipt encodes the EBITDA rework reduction claim
(\$1.25M), carrying fitness $0.982$ and precision $0.945$ from the conformance authority
query. The \texttt{rec\_wc\_ar\_002.json} receipt encodes the working capital and accounts
receivable velocity claim (\$1.37M), carrying fitness $0.978$ with replication verdict
\texttt{PASSED}. The \texttt{rec\_risk\_sla\_003.json} receipt encodes the SLA penalty
exposure cap (\$450K), with a late delivery rate of $2.14\%$ and log coverage of $98.4\%$.
The \texttt{rec\_risk\_compliance\_004.json} receipt encodes the GRC compliance leakage
claim (\$0.00 leakage, fitness $1.0$, zero violations). The
\texttt{rec\_residual\_standard\_005.json} receipt encodes the residual process
standardization claim ($97.5\%$ conformance, drift status \texttt{STABLE}).

The module generation receipt layer contains five Markdown documents. The
\texttt{wasm4pm\_conformance\_authority\_generation.md} is the most structurally significant,
documenting version \texttt{v30.1.2} of the 664-line Rust conformance authority module with
Blake3 hash \texttt{75508e6e9e7acccb679e5c8be35ddd2adc9c4732f1fc331047d755703c9107c3} and
a record of $8/8$ unit tests passing. The \texttt{wasm4pm\_mining\_generation.md},
\texttt{wasm4pm\_replay\_generation.md}, and \texttt{wasm4pm\_lifecycle\_generation.md}
document the mining, replay, and lifecycle module manufacturing pipelines respectively.
The \texttt{blue\_river\_generation.md} documents the Blue River Dam orchestrator receipt
with status \texttt{ORCHESTRATOR\_ALIVE} and $5/5$ unit tests passing.

The governance layer contains three documents: the \texttt{RECEIPT\_REGISTRY.md} enumerating
seven canonical receipt types, the \texttt{ma\_deck\_rendering\_authority\_assessment.md}
confirming the seven-step board presentation pipeline, and the related checkpoint files
in \texttt{checkpoints/}.

\subsection{Commands}
\label{sec:receipts-implementation-commands}

The manufacturing pipeline for M\&A receipts is driven by the \texttt{wasm4pm} Rust binary
compiled from the conformance authority module. The conformance authority exposes a
\texttt{ConformanceAdmissionGate} callable from the module's public interface. The
gate evaluates a \texttt{ConformanceVerdicts} struct and either passes the receipt
forward for signing or emits a hard-reject verdict.

Module integrity is verified through the \texttt{cargo check} command, which confirmed
zero compilation errors across the conformance authority module (8 warnings unrelated to
conformance authority logic). The full test suite is exercised via \texttt{cargo test},
which executes the seven named unit tests: \texttt{test\_admission\_gate\_pass},
\texttt{test\_token\_fitness\_perfect}, \texttt{test\_admission\_gate\_conditional\_with\_override},
\texttt{test\_token\_fitness\_degraded}, \texttt{test\_admission\_gate\_hard\_reject},
\texttt{test\_verdict\_display}, and \texttt{test\_reachability\_heuristic}.

The Blue River Dam orchestrator, documented in \texttt{blue\_river\_generation.md}, is compiled
and tested with the same \texttt{cargo check} and \texttt{cargo test} workflow, producing
$5/5$ test passes for \texttt{test\_governance\_hierarchy}, \texttt{test\_lifecycle\_state\_machine},
\texttt{test\_quality\_gate\_soundness}, \texttt{test\_conformance\_metric}, and
\texttt{test\_mape\_k\_knowledge\_persistence}.

\subsection{Data and Ontology Surfaces}
\label{sec:receipts-implementation-data}

The data surface of the receipts corpus is defined by the JSON schema of the M\&A claim
receipts and the Markdown schema of the generation receipts. The JSON receipts follow
JCS canonical form (RFC 8785) to ensure deterministic serialization for BLAKE3 hashing
and Ed25519 signing. The \texttt{target\_log\_hash} field in each JSON receipt is a
BLAKE3 hash of the XES or OCEL event log file from which the conformance measurement
was derived. The \texttt{process\_model\_hash} field is a BLAKE3 hash of the process
model (typically a Petri net or process tree) used as the reference for conformance checking.

The receipt registry defines an ontology of seven canonical receipt types, each with
a named witness list and a gap inventory. This registry is the authority for what
constitutes a receipted fact in the research program versus an asserted claim. The
registry distinguishes between \texttt{PAPER\_CANON} receipts (anchoring paper-to-type-law
mappings), \texttt{PM4PY\_ORACLE} receipts (anchoring pm4py capability assessments),
\texttt{WASM4PM\_GAP} receipts (anchoring gap findings), \texttt{LIFECYCLE} receipts
(anchoring state machine transitions), \texttt{MA} receipts (the five financial claim
receipts), \texttt{STANDARDS} receipts (anchoring standards compliance maps), and
\texttt{ADVERSARIAL} receipts (anchoring hostile test case results).

\subsection{Templates and Rendered Artifacts}
\label{sec:receipts-implementation-templates}

The board presentation layer of the receipts corpus is implemented through two Tera2
templates: \texttt{ma-deck.tera} (202 lines) for the executive PowerPoint presentation
and \texttt{ma-diligence.tera} (283 lines) for the technical due diligence workbook.
These templates consume the signed JSON receipts as data sources, binding each slide
and worksheet cell to a specific receipt field. The rendering pipeline uses the Rust
libraries \texttt{pptx-rs} for PowerPoint artifact manufacture and \texttt{calamine}
for Excel workbook rendering.

The \texttt{ma\_deck\_rendering\_authority\_assessment.md} confirms that all seven
steps of this board presentation pipeline are either \texttt{VERIFIED} (template exists,
data binding confirmed) or \texttt{EXECUTABLE} (pipeline ready, awaiting final execution
trigger). As of the \texttt{PROCESS\_INTELLIGENCE\_ALIVE\_001} checkpoint, the pipeline
is assessed as complete and ready for execution. The final rendered artifacts---the
\texttt{.pptx} board deck and the \texttt{.xlsx} due diligence workbook---are identified
as \texttt{FUTURE WORK}: they are designated in the pipeline but have not yet been
materialized as files in the repository. This is documented as an open gap in
Section~\ref{sec:receipts-questions-gaps}.
